% GENERATED FILE. Edit canonical lessons, then run npm run generate:books.
% canonical-source-hash: fnv1a64:736c142d5ce0bc55
% canonical-lessons: KA-C61-curd, KA-C61-ghee, KA-C61-oil, KA-C61-jaggery, KA-C61-sugar

\chapter{In the Kitchen}
\label{ch:ka-kitchen}
\hlchaptermodality{61}

\begin{chapteropening}
\textbf{By the end of this chapter:} \emph{I can name five things in a Kannada kitchen, from the curd to the sugar.}

Complete KA-C61-sugar and say all five words in order.
\end{chapteropening}

\section[mosaru]{\kn{ಮೊಸರು} (mosaru) — curd}

\label{lesson:KA-C61-curd}



\begin{quote}
Before the new run: what did \emph{kāge} mean, and what did \emph{hakki} mean?
\end{quote}

\subsection*{You'll want to know: \kn{ಮೊಸರು}}
\textbf{\kn{ಮೊಸರು}} (\emph{mosaru}) — \textquotedblleft{}curd\textquotedblright{}.

\kn{ಮೊಸರು} is what \kn{ಹಾಲು} becomes when it is left overnight with a spoonful of yesterday's \kn{ಮೊಸರು} in it. The \kn{ಹಸು} from the last chapter is the front of that chain and this is the end of it.

It closes a meal rather than opening one. \kn{ಮೊಸರು} over \kn{ಅನ್ನ} is the last thing on the plate, and a Kannada speaker will tell you the meal is not finished without it.

\begin{grammarlens}[title={What you've built}]
The first of five things in a Kannada kitchen.
\end{grammarlens}

\subsection*{Your turn}
Say these aloud:

\begin{itemize}
\raggedright
  \item \emph{mosaru}
  \item it once more, slowly
  \item \emph{hālu}, then \emph{mosaru}, and say which one waits overnight
\end{itemize}

\begin{tcolorbox}[breakable,colback=teal!4,colframe=teal!35!black,title={Before you move on}]
What does \kn{ಮೊಸರು} mean? (\textquotedblleft{}curd\textquotedblright{}.) Say it once more.
\end{tcolorbox}
\section[tuppa]{\kn{ತುಪ್ಪ} (tuppa) — ghee}

\label{lesson:KA-C61-ghee}



\begin{quote}
Before the new one: what did \emph{mosaru} mean?
\end{quote}

\subsection*{You'll want to know: \kn{ತುಪ್ಪ}}
\textbf{\kn{ತುಪ್ಪ}} (\emph{tuppa}) — \textquotedblleft{}ghee\textquotedblright{}.

\kn{ತುಪ್ಪ} is the other thing \kn{ಹಾಲು} turns into, further along the same road: cream taken off, churned, then cooked until only the fat is left. Inherited, and shared with the sisters.

A spoonful goes on hot \kn{ಅನ್ನ}, and a \kn{ದೀಪ} will burn on it as readily as on any oil. It is the one item in this chapter that is food and light at once.

\begin{grammarlens}[title={What you've built}]
Two.
\end{grammarlens}

\subsection*{Your turn}
Say these aloud:

\begin{itemize}
\raggedright
  \item \emph{tuppa}
  \item it once more, slowly
  \item \emph{mosaru}, then \emph{tuppa}, and say which one is poured
\end{itemize}

\begin{tcolorbox}[breakable,colback=teal!4,colframe=teal!35!black,title={Before you move on}]
What does \kn{ತುಪ್ಪ} mean? (\textquotedblleft{}ghee\textquotedblright{}.) Say it once more.
\end{tcolorbox}
\section[eṇṇe]{\kn{ಎಣ್ಣೆ} (eṇṇe) — oil}

\label{lesson:KA-C61-oil}



\begin{quote}
Before the new one: what did \emph{tuppa} mean?
\end{quote}

\subsection*{You'll want to know: \kn{ಎಣ್ಣೆ}}
\textbf{\kn{ಎಣ್ಣೆ}} (\emph{eṇṇe}) — \textquotedblleft{}oil\textquotedblright{}.

\kn{ಎಣ್ಣೆ} is Kannada's own oil-word, and it owes \kn{ತುಪ್ಪ} nothing at all. \kn{ತುಪ್ಪ} comes out of an animal; \kn{ಎಣ್ಣೆ} is pressed out of a \kn{ಬೀಜ}, and the two sit in different jars on the same shelf.

There is a small joke in this for you. The friend-word \kn{ಸ್ನೇಹಿತ} was built on a Sanskrit word whose first sense was oil — smoothness, the thing that makes two surfaces run easily on each other. That is the borrowed, figurative oil. This is the one in the jar.

\begin{grammarlens}[title={What you've built}]
Three.
\end{grammarlens}

\subsection*{Your turn}
Say these aloud:

\begin{itemize}
\raggedright
  \item \emph{eṇṇe}
  \item it once more, slowly
  \item \emph{tuppa}, then \emph{eṇṇe}, and say which one starts as a seed
\end{itemize}

\begin{tcolorbox}[breakable,colback=teal!4,colframe=teal!35!black,title={Before you move on}]
What does \kn{ಎಣ್ಣೆ} mean? (\textquotedblleft{}oil\textquotedblright{}.) Say it once more.
\end{tcolorbox}
\section[bella]{\kn{ಬೆಲ್ಲ} (bella) — jaggery}

\label{lesson:KA-C61-jaggery}



\begin{quote}
Before the new one: what did \emph{eṇṇe} mean?
\end{quote}

\subsection*{You'll want to know: \kn{ಬೆಲ್ಲ}}
\textbf{\kn{ಬೆಲ್ಲ}} (\emph{bella}) — \textquotedblleft{}jaggery\textquotedblright{}.

\kn{ಬೆಲ್ಲ} is cane juice boiled down and set in a lump, brown and sticky and sold by weight rather than by the spoon. Kannada's own word, and the older of the two sweet things in this chapter.

South Karnataka cooking puts a piece of it into food that is not a sweet at all — into the same dish as the \emph{uppu} you already have, deliberately, so that neither taste wins. A cook who leaves it out is making somebody else's version of the dish.

\begin{grammarlens}[title={What you've built}]
Four.
\end{grammarlens}

\subsection*{Your turn}
Say these aloud:

\begin{itemize}
\raggedright
  \item \emph{bella}
  \item it once more, slowly
  \item \emph{bella}, then \emph{uppu}, and say which two go into the same pot
\end{itemize}

\begin{tcolorbox}[breakable,colback=teal!4,colframe=teal!35!black,title={Before you move on}]
What does \kn{ಬೆಲ್ಲ} mean? (\textquotedblleft{}jaggery\textquotedblright{}.) Say it once more.
\end{tcolorbox}
\section[sakkare]{\kn{ಸಕ್ಕರೆ} (sakkare) — sugar}

\label{lesson:KA-C61-sugar}



\begin{quote}
Before the new one: what did \emph{bella} mean?
\end{quote}

\subsection*{You'll want to know: \kn{ಸಕ್ಕರೆ}}
\textbf{\kn{ಸಕ್ಕರೆ}} (\emph{sakkare}) — \textquotedblleft{}sugar\textquotedblright{}.

\kn{ಸಕ್ಕರೆ} is the refined one, white and loose, and it is the traveller of this chapter. It goes back to a Sanskrit word meaning gravel or grit — a fair description of what sugar looks like beside a lump of \kn{ಬೆಲ್ಲ}.

The coffee-word and the tea-word each reached Kannada after a long journey inward. This word went the other way. It left India, crossed Arabic and Latin, turned into the English \emph{sugar}, and Kannada kept the original at home the whole time.

\begin{grammarlens}[title={What you've built}]
Five: \kn{ಮೊಸರು}, \kn{ತುಪ್ಪ}, \kn{ಎಣ್ಣೆ}, \kn{ಬೆಲ್ಲ}, \kn{ಸಕ್ಕರೆ}. Enough to say what is in the food you are given.
\end{grammarlens}

\subsection*{Your turn}
Say these aloud:

\begin{itemize}
\raggedright
  \item \emph{sakkare}
  \item it once more, slowly
  \item all five in order, then give your name to the person cooking
\end{itemize}

\begin{tcolorbox}[breakable,colback=teal!4,colframe=teal!35!black,title={Before you move on}]
What does \kn{ಸಕ್ಕರೆ} mean? (\textquotedblleft{}sugar\textquotedblright{}.) Say it once more.
\end{tcolorbox}
